%% ============================================================
\section{Chain-of-Evidence: A Standard for Research Verifiability}
\label{sec:coe}
%% ============================================================



\begin{quote}
\emph{\textbf{Principle:} Every claim produced by a research system must be traceable, through a recorded chain of supporting claims and evidence, to a grounding source.}
\end{quote}

\noindent A credible research claim must be backed by verifiable evidence.
Without this requirement, the same system that produces a plausible-sounding paper can also produce fabricated citations, hallucinated numbers, and descriptions of experiments that never happened.
Just as a database that violates ACID may return plausible-looking query results even as it silently corrupts data---a transfer debits one account but never credits another, yet both balances look valid---a research system that violates CoE may produce plausible-looking papers whose claims cannot be traced to evidence---the paper reads well, but the scores do not reproduce.
ACID does not prescribe \emph{how} to build a database. It prescribes what properties the database must have.
CoE plays the same role for research artifacts.


We define four primary claim types, each with a required evidence chain shape. The taxonomy is not exhaustive but covers the claim types that are tractably verifiable with current tools---other types (e.g., qualitative observations, theoretical properties) require domain expertise or subjective judgment that is harder to automate.
\textbf{\emph{Citation claims}} (e.g., ``Smith et al.\ showed X'') require that the cited work exists in a scholarly database and that its content is consistent with how it is described in the paper.
\textbf{\emph{Numerical claims}} (e.g., ``achieves 87.3\% on Prism'') must trace from the reported value to a recorded output (e.g., an execution log, experimental measurement, or simulation result).
\textbf{\emph{Methodological claims}} (e.g., ``we use a 3-layer MLP'') must resolve from the method description to the corresponding implementation.
\textbf{\emph{Conclusion claims}} (e.g., ``outperforms baseline by 5\%'') must derive from supporting claims---numerical, methodological, or both---through verifiable reasoning.
CoE is deliberately architecture-agnostic: it defines what properties a verifiable artifact should have, not how the system should construct one.
The standard is also author-agnostic---the same evidence chains are required whether a paper is human- or machine-authored---but we focus on autonomous systems because their failure modes are systematic and rapidly growing in scale.

In the following sections, we describe \sys{}, an autonomous research system designed to satisfy CoE by construction (\S\ref{sec:system}), and \coeaudit{}, a post-hoc audit that measures how well any system's artifacts meet the standard through four integrity checks (\S\ref{sec:coeaudit}).
